\paragraph*{Three models.} Three logistic regression models score each scan: one on the images, one on the clinical
history and one on both. The imaging model uses the 27 measurements and is never given dates. The clinical-history
model uses the treatment course and the patient's characteristics, and the combined model uses both input sets.
PRISM shows all three, so the record shows which kind of evidence a score follows. Each model is a logistic regression
with an L2 (ridge) penalty \citep{pedregosa2011}, chosen because it is simple and uses few inputs. Deep networks tuned
on MU did not improve on it (Table~\ref{tab:2}). Inputs are standardised, and the penalty is chosen from six values of
the inverse strength $C$ (0.01 to 3) by threefold cross validation grouped by patient, minimising log loss.

LUMIERE records only the days since surgery, and UCSF-ALPTDG the days since surgery and the start of radiotherapy,
so the full
clinical-history model cannot be computed there. For these centres PRISM uses reduced models, fitted on MU with the
inputs that every cohort records. The reduced clinical-history model takes age, sex, IDH and MGMT status, days since
surgery and a flag for the first weeks after radiotherapy. At LUMIERE the flag is the rater's mark for scans within
three months of radiotherapy, and at UCSF-ALPTDG it is approximated from the start of radiotherapy, assuming a six
week course. The reduced combined model adds the enhancing, necrosis (with cavity) and FLAIR volumes, their change
since the previous scan, and the interval since that scan. LUMIERE's masks have three compartments, so its imaging
model is fitted on MU with those three volumes and their change. Because the imaging model needs only scans and masks,
PRISM gives an imaging score at a centre that records no treatment dates.

% Sources: prism/audit/conformal.py; METHODS_PLAN.md M2; RELIABILITY_CARD.md; REVISION_DATA.md R2; refs
% angelopoulos2023, vovk2022book, romano2020, barber2021limits, johansson2015, ding2023, mehrtens2023, mitchell2019.
